% Scope (SPEC table of contents): The interpretability debate; the evaluation taxonomy; faithfulness vs plausibility defined; the book's thesis stated. Papers 13, 14, 15.
\chapter{Giải thích mô hình để làm gì?}
\label{ch:giai-thich-de-lam-gi}

Một mô hình có thể dự đoán đúng mà vẫn đặt ra câu hỏi chưa được trả lời: nó
dựa vào đâu để đi tới dự đoán ấy, và lý do ta đọc được có thực sự là lý do mô
hình đã dùng hay không? Hai câu hỏi này không giống nhau. Câu đầu nói về việc
con người cần biết gì trước khi giao một quyết định cho mô hình. Câu sau nói về
độ đúng của chính lời giải thích.

Chương này tách các câu hỏi đó ra trước khi bàn đến LIME, SHAP hay attribution.
Ta sẽ gọi thuộc tính cho biết lời giải thích có bám vào cơ chế ra quyết định của
mô hình là \tn{độ trung thực}{faithfulness}. Một lời giải thích nghe hợp lý
với người đọc có thể thiếu thuộc tính ấy. Phân biệt này là tiền đề cho toàn bộ
phần phê phán ở nửa sau cuốn sách.

\section{Khi dự đoán chưa đủ}
\label{sec:ch01-khi-du-doan-chua-du}

\index{khả năng diễn giải|(}%
Một hệ thống học máy thường được đánh giá bằng đầu ra: dự đoán có đúng trên
dữ liệu kiểm tra không, lỗi có nằm trong mức chấp nhận được không. Cách đánh
giá này trả lời một câu hỏi hẹp về hiệu năng, nhưng không tự trả lời liệu hệ
thống có đang dùng một quy luật mà ta chấp nhận được, có còn phù hợp khi môi
trường thay đổi, hay có đẩy rủi ro sang một nhóm người mà dữ liệu kiểm tra
không làm lộ ra.

Nhu cầu về \tn{khả năng diễn giải}{interpretability} nảy sinh từ chính khoảng
trống đó. Tuy vậy, Lipton chỉ ra rằng từ này gộp nhiều mục đích khác nhau:
tạo niềm tin, tìm quan hệ nhân quả, chuyển giao sang tình huống mới, lấy thêm
thông tin cho người ra quyết định, và bảo đảm tính công bằng~\autocite{p13mythos}.
Gộp các mục đích ấy dưới một nhãn khiến một phương
pháp có thể được khen là \enquote{dễ hiểu} mà không ai nói rõ ai sẽ dùng nó,
để làm việc gì, và bằng tiêu chuẩn nào.

Xét một mô hình dự báo nguy cơ cho bệnh nhân. Một điểm số chính xác trên tập
kiểm tra chưa cho biết mô hình đang dùng tín hiệu bệnh lý hay đang tận dụng
một dấu vết của quy trình điều trị cũ. Bác sĩ có thể cần lời giải thích để
kiểm tra khả năng thứ hai, nhà nghiên cứu có thể cần nó để nêu một giả thuyết
về bệnh, còn người quản lý có thể cần nó để xem một tiêu chí công bằng đã bị
bỏ sót hay chưa. Như vậy, cùng một lời giải thích nhưng ba người dùng cần nó
cho ba việc kiểm tra khác nhau.

\index{hộp đen}%
Vì vậy, ta không bắt đầu bằng câu hỏi mô hình nào là \tn{hộp đen}{black box},
mà bằng câu
hỏi: điều gì còn nằm ngoài hàm mục tiêu và phép đo hiện có? Doshi-Velez và
Kim gọi tình trạng này là sự không đầy đủ của đặc tả bài toán: điều ta quan
tâm không thể liệt kê trọn vẹn để tối ưu hoặc để kiểm thử trực tiếp~\autocite{p14rigorous}.
Lời giải thích không tự lấp khoảng trống đó; nó chỉ đưa thêm bằng chứng để
con người kiểm tra khoảng trống ấy.

\index{khả năng diễn giải|)}%

\section{Giải thích là gì?}
\label{sec:ch01-giai-thich-la-gi}

Trong cuốn sách này, một giải thích là thông tin được trình bày để con người
hiểu một khía cạnh của hành vi mô hình. Định nghĩa này cố ý chưa hứa rằng thông
tin ấy mô tả đúng cơ chế bên trong. Hứa như vậy quá sớm sẽ trộn lẫn hai việc:
trình bày cho người đọc và mô tả điều mô hình đã làm.

\begin{definition}[\tn{Khả năng diễn giải}{interpretability}]
\label{def:ch01-kha-nang-dien-giai}
Khả năng diễn giải là năng lực trình bày hành vi của một hệ thống học máy bằng
những thuật ngữ con người có thể hiểu và dùng cho một mục đích đã nêu.
\end{definition}

Định nghĩa~\ref{def:ch01-kha-nang-dien-giai} có hai phần đáng giữ lại. Thứ
nhất, người nhận giải thích là một phần của bài toán. Bản đồ nhiệt có thể hữu
ích cho chuyên gia ảnh y khoa nhưng vô dụng trong một quy trình khi người vận
hành phải chọn giữa các phương án hành động. Thứ hai, tính hữu ích không đồng
nghĩa với tính đúng về cơ chế.

Một hướng tiếp cận đặt \tn{minh bạch}{transparency} vào chính mô hình. Với
một mô hình tuyến tính nhỏ, ta có thể lần theo biến đầu vào, tham số và phép
tính dẫn đến một dự đoán. Hướng khác tạo \tn{giải thích hậu nghiệm}{post-hoc
explanation} sau khi mô hình đã được huấn luyện: một danh sách đặc trưng, một
bản đồ nhiệt, hoặc một mô hình thay thế cục bộ. Lipton đặt hai hướng này cạnh
nhau vì chúng trả lời những câu hỏi khác nhau~\autocite{p13mythos}.

Không nên biến sự phân biệt này thành một thứ bậc đơn giản. Mô hình minh bạch
cũng có thể quá lớn để người dùng kiểm tra toàn bộ, còn một giải thích hậu
nghiệm có thể giúp phát hiện một lỗi cụ thể. Điểm cần giữ là nguồn gốc của
lời giải thích. Khi nó chỉ xấp xỉ hành vi của mô hình, ta cần một cách kiểm tra
riêng cho quan hệ giữa xấp xỉ và cơ chế.

\section{Đánh giá một lời giải thích}
\label{sec:ch01-danh-gia-giai-thich}

Mục~\ref{sec:ch01-khi-du-doan-chua-du} đã cho thấy mỗi mục đích đòi một việc
kiểm tra riêng. Nếu mục đích quyết định thước đo thì không thể đánh giá mọi
lời giải thích bằng một con số hoặc một đại lượng thay thế (proxy) mang tính
cấu trúc. Doshi-Velez và Kim đề xuất ba cách đặt bài toán đánh giá, khác nhau
ở người tham gia và nhiệm vụ được dùng~\autocite{p14rigorous}.
Hình~\ref{fig:ch01-ba-cach-danh-gia} đặt ba cách ấy trên cùng một trục.

\begin{figure}[htbp]
  \centering
  \begin{tikzpicture}[
  box/.style={draw=black!65, rounded corners=2pt, align=center,
    minimum width=4.4cm, minimum height=1.1cm, fill=black!5},
  data/.style={draw=black!65, rounded corners=2pt, align=center,
    minimum width=4.4cm, minimum height=1.1cm, fill=black!15},
  arrow/.style={-{Stealth[length=2mm]}, draw=black!65, thick}
]
  \node[data] (application) {Đánh giá trong ứng dụng};
  \node[box, below=18mm of application] (human) {Đánh giá với người dùng};
  \node[box, below=18mm of human] (functional) {Đánh giá theo chức năng};

  \draw[arrow] (human) -- (application);
  \draw[arrow] (functional) -- (human);

  \node[align=left, anchor=west] at ($(human.north)!0.5!(application.south) + (8mm,0)$)
    {vay giả định rằng nhiệm vụ rút gọn\\giữ được phần cốt lõi của nhiệm vụ thật};
  \node[align=left, anchor=west] at ($(functional.north)!0.5!(human.south) + (8mm,0)$)
    {vay giả định rằng đại lượng thay thế\\đã được kiểm chứng với người dùng};
\end{tikzpicture}

  \caption{Ba cách đánh giá khác nhau ở bằng chứng mà chúng đòi hỏi. Đánh giá
  trong ứng dụng dùng người dùng và nhiệm vụ thực nên gần mục tiêu triển khai
  nhất. Đánh giá với người dùng rút gọn nhiệm vụ để cô lập một giả thuyết.
  Đánh giá theo chức năng dùng proxy, rẻ hơn nhưng phải biện minh vì sao proxy
  có liên hệ với điều người dùng cần. Sự tiện lợi của cách rẻ hơn không chuyển
  thành bằng chứng cho cách gần mục tiêu hơn.}
  \label{fig:ch01-ba-cach-danh-gia}
\end{figure}

\paragraph{Đánh giá trong ứng dụng.} Người dùng thật thực hiện nhiệm vụ thật;
ví dụ, một bác sĩ dùng lời giải thích để phát hiện khi mô hình dự báo dựa vào
một tín hiệu không phù hợp. Đây là cách đánh giá gần với mục đích triển khai
nhất, nhưng cũng tốn kém và khó kiểm soát các yếu tố nhiễu.

\paragraph{Đánh giá với người dùng.} Người tham gia vẫn là con người, nhưng
nhiệm vụ được giản lược; chẳng hạn, ta có thể yêu cầu họ chọn lời giải thích
nào giúp họ dự đoán hành vi mô hình tốt hơn. Thiết kế này không tái tạo toàn
bộ công việc thực tế, nhưng cho phép kiểm tra một thuộc tính xác định trong
cách trình bày lời giải thích mà không phải dựng cả ứng dụng.

\paragraph{Đánh giá theo chức năng.} Cách thứ ba không cần người tham gia:
nghiên cứu đo một proxy như độ thưa của mô hình, độ ngắn của quy tắc, hay khả
năng dự đoán một nhãn thay thế. Cách làm này phù hợp khi proxy đã có cơ sở,
nhưng proxy không phải bản thân mục tiêu. Nghiên cứu so sánh các mô hình dễ
diễn giải gần đây cũng phải chọn proxy cấu trúc để có thể so sánh trên quy mô
lớn, và thừa nhận giới hạn của lựa chọn ấy~\autocite{p15microscope}.

Như vậy, ba cách trên không phải là ba nhãn chất lượng cho cùng một đại
lượng, mà là ba câu trả lời cho ba câu hỏi thực nghiệm khác nhau. Một kết quả
tốt ở cách đánh giá theo chức năng có thể là lý do để làm nghiên cứu tiếp,
không phải giấy chứng nhận rằng người dùng sẽ hiểu đúng hoặc rằng lời giải
thích đã trung thực.

\section{Độ trung thực và tính hợp lý}
\label{sec:ch01-do-trung-thuc-va-tinh-hop-ly}

\index{tính hợp lý}%
Dù được đánh giá theo cách nào trong ba cách ở mục trước, một lời giải thích
vẫn có thể thuyết phục được người đọc chỉ vì nó ngắn, quen thuộc và phù hợp
với trực giác. Những phẩm chất đó quan trọng cho giao tiếp, nhưng chưa nói gì
về việc mô hình đã tính toán thế nào. Ta gọi sự phù hợp với đánh giá của
người đọc là \tn{tính hợp lý}{plausibility}; đây là một thuộc tính của lời
giải thích xét trong quan hệ với người nhận.

\index{độ trung thực}%
\begin{definition}[Độ trung thực]
\label{def:ch01-do-trung-thuc}
Lời giải thích có \tn{độ trung thực}{faithfulness} đối với một hành vi và một
tập can thiệp đã xác định khi những đặc trưng mà nó xếp là quan trọng đúng là
những đặc trưng mà hành vi ấy phụ thuộc vào dưới các can thiệp thuộc tập đó,
và những đặc trưng mà nó xếp là không quan trọng thì hành vi ấy không phụ
thuộc vào.
\end{definition}

Chữ phụ thuộc trong định nghĩa~\ref{def:ch01-do-trung-thuc} mang nghĩa của
can thiệp: hành vi phụ thuộc vào một đặc trưng khi một can thiệp được phép lên
đặc trưng ấy làm hành vi thay đổi. Cách phát biểu này là cách sách viết lại
định nghĩa mà phần lớn tài liệu dùng, của Jacovi và Goldberg, theo đó lời
giải thích trung thực là lời giải thích phản ánh đúng quá trình suy luận đằng
sau tiên đoán của mô hình; chương~\ref{ch:chi-so-khong-do-do-trung-thuc} sẽ
dẫn lại định nghĩa ấy qua bài báo neo của cuốn sách~\autocite{p21metrics}.
Sách thêm hai vế hành vi và tập can thiệp vì quá trình suy luận của một mô
hình không quan sát trực tiếp được, còn phản ứng của một hành vi trước một can
thiệp thì đo được.

Định nghĩa~\ref{def:ch01-do-trung-thuc} vì thế buộc ta nêu rõ hai đối tượng
trước khi đánh giá. Thứ nhất, hành vi nào đang được giải thích: một dự đoán, một
lớp đầu ra, hay một vùng đầu vào? Thứ hai, can thiệp nào được phép: che một
đặc trưng, thay nó bằng một giá trị nền, hay thay đổi một khái niệm ở mức
cao? Khi hai công trình trả lời khác nhau cho hai câu này, kết quả của chúng
không nhất thiết mâu thuẫn, vì chúng có thể đã đo hai quan hệ khác nhau.

Tính hợp lý và độ trung thực thường đi cùng nhau trong những ví dụ thuận lợi,
vì vậy chúng dễ bị lẫn; nhưng chúng tách ra ngay khi lời giải thích kể một
câu chuyện quen thuộc mà mô hình không hề dựa vào. Một phương pháp hậu nghiệm
có thể cho lời giải thích mạch lạc trong khi bỏ qua đặc trưng thực sự chi
phối dự đoán. Billa và cộng sự nêu cùng điểm ấy ở dạng hẹp hơn: vì phương pháp
hậu nghiệm chỉ là xấp xỉ, nó có thể không phản ánh đúng cách mô hình đã ra
quyết định~\autocite{p15microscope}.

Từ đây, mỗi lần cuốn sách dùng chữ attribution, câu hỏi đi kèm sẽ là:
attribution đang được xem như một cách trình bày cho người đọc, hay như bằng
chứng về cơ chế? Phần~II của cuốn sách khảo sát các phương pháp trả lời câu
hỏi này bằng những giả định khác nhau.

\section{Luận đề của cuốn sách}
\label{sec:ch01-luan-de-cua-sach}

Luận đề của cuốn sách không phải là mọi giải thích đều vô ích, cũng không
phải là mô hình phức tạp luôn phải bị thay bằng mô hình đơn giản. Luận đề hẹp
hơn thế: một thước đo được gọi là \tn{độ trung thực}{faithfulness} phải được
kiểm chứng như một dụng cụ đo. Nếu chưa biết điểm số của nó thay đổi ra sao
khi cơ chế thật thay đổi, hoặc nó có phân biệt được lời giải thích đúng cơ
chế với lời giải thích chỉ nghe hợp lý hay không, thì ta chưa có lý do để xem
điểm số ấy là bằng chứng về độ trung thực.

Các chương kế tiếp đi theo thứ tự cần thiết cho lập luận đó.
Chương~\ref{ch:mo-hinh-can-giai-thich} thu gọn nền tảng về các mô hình mà
những lời giải thích sau này sẽ nhắm tới. Các chương về LIME, SHAP và
attribution theo gradient trình bày cơ chế, giả định và đầu ra của từng họ
phương pháp. Sau khi cơ chế, giả định và đầu ra của từng họ đã rõ,
chương~\ref{ch:do-do-trung-thuc} phân loại các họ chỉ số độ trung thực, còn
chương~\ref{ch:chi-so-khong-do-do-trung-thuc} kiểm tra trực diện tuyên bố
rằng các chỉ số ấy đo được độ trung thực.

Trình tự này cũng nhằm tránh một cách đọc sai phổ biến: khi một bản đồ nhiệt
hoặc danh sách đặc trưng trông thuyết phục, ta dễ nhảy thẳng từ hình thức của
lời giải thích sang kết luận về mô hình. Cuốn sách giữ hai bước tách biệt:
trước hết mô tả phương pháp đang tạo ra cái gì, sau đó mới hỏi bằng chứng nào
cho thấy nó bám vào hành vi mô hình. Câu hỏi sau khó hơn, và chính nó quyết
định ta được phép đặt bao nhiêu niềm tin vào lời giải thích.

\begin{tomtat}
\begin{itemize}
  \item Hiệu năng dự đoán không tự trả lời các câu hỏi về cơ chế, mục tiêu còn
  thiếu trong đặc tả, hay cách con người sẽ dùng mô hình.
  \item \tn{Khả năng diễn giải}{interpretability} là một mục tiêu gắn với
  người nhận và công việc; \tn{minh bạch}{transparency} và \tn{giải thích hậu
  nghiệm}{post-hoc explanation} có nguồn gốc khác nhau.
  \item Đánh giá trong ứng dụng, đánh giá với người dùng và đánh giá theo chức
  năng dùng các loại bằng chứng khác nhau.
  \item \tn{Tính hợp lý}{plausibility} với người đọc không chứng minh
  \tn{độ trung thực}{faithfulness} với cơ chế mô hình.
  \item Một chỉ số độ trung thực là dụng cụ đo và cần được kiểm chứng như một
  dụng cụ đo.
\end{itemize}
\end{tomtat}

\cauhoi
\begin{cauhoids}
  \item Nêu một tình huống trong đó mô hình đạt hiệu năng tốt nhưng người dùng
  vẫn cần một lời giải thích. Xác định điều gì còn thiếu trong đặc tả bài toán.
  \item Chọn một lời giải thích hậu nghiệm quen thuộc. Theo
  định nghĩa~\ref{def:ch01-do-trung-thuc}, hãy nêu hành vi và loại can thiệp
  cần xác định trước khi có thể hỏi nó có trung thực hay không.
  \item Một đánh giá theo chức năng cho thấy mô hình có ít quy tắc hơn mô hình
  khác. Kết quả đó hỗ trợ được kết luận nào, và không hỗ trợ được kết luận nào
  về việc người dùng có hiểu mô hình hay không?
\end{cauhoids}

